\documentclass[a4paper,12pt]{ctexart}
\usepackage{xcolor,graphicx,amsmath}
\usepackage[top=1.8cm, bottom=1.8cm, left=1.8cm, right=1.8cm]{geometry}
\usepackage{array,longtable,hyperref}
\hypersetup{colorlinks=true,linkcolor=blue!70!black}
\linespread{1.35}
\definecolor{defblue}{RGB}{0,80,160}\definecolor{thinkorange}{RGB}{230,120,20}
\definecolor{formulared}{RGB}{180,30,50}\definecolor{funboxbg}{RGB}{245,248,252}
\definecolor{examplebg}{RGB}{255,250,240}\definecolor{practicegreen}{RGB}{40,130,70}
\newenvironment{thinkbox}{\vspace{0.3cm}\begin{center}\begin{minipage}{0.88\textwidth}\colorbox{funboxbg}{\begin{minipage}{0.94\textwidth}\vspace{0.15cm}\textbf{\textcolor{thinkorange}{【 想一想 】}}}{\vspace{0.15cm}\end{minipage}}\end{minipage}\end{center}\vspace{0.3cm}}
\newenvironment{definition}[1]{\vspace{0.2cm}\begin{center}\begin{minipage}{0.88\textwidth}\fbox{\begin{minipage}{0.92\textwidth}\vspace{0.1cm}\textbf{\textcolor{defblue}{【 #1 】}}}{\vspace{0.1cm}\end{minipage}}\end{minipage}\end{center}\vspace{0.2cm}}
\newenvironment{practice}{\vspace{0.3cm}\begin{center}\begin{minipage}{0.88\textwidth}\colorbox{funboxbg}{\begin{minipage}{0.94\textwidth}\vspace{0.15cm}\textbf{\textcolor{practicegreen}{【 练一练 】}}}{\vspace{0.15cm}\end{minipage}}\end{minipage}\end{center}\vspace{0.2cm}}
\newenvironment{figplaceholder}[1]{\vspace{0.2cm}\begin{center}\fbox{\begin{minipage}{0.82\textwidth}\centering\textcolor{blue!70!black}{\textbf{【插图位置】}}\\[0.2cm]#1\end{minipage}}\end{center}\vspace{0.2cm}}
\title{\textcolor{blue!70!black}{\Huge\textbf{生物1 · 分子与细胞}}}
\author{}\date{}
\begin{document}
\maketitle\thispagestyle{empty}
\section*{\textcolor{blue!70!black}{致同学}}
欢迎来到生物学的世界。如果说物理研究的是"力"、化学研究的是"物质"，那么生物学研究的就是"活的东西"。这一本将从分子和细胞的最底层开始——看看构成你身体的60万亿个细胞是怎么工作的。

\newpage
\section{\textcolor{orange!80!black}{第一课\quad 生物的特征——什么是"活着"}}

\subsection*{什么算"活"？}
一块石头是活的吗？一颗种子是活的吗？一个电脑程序是活的吗？

生物和非生物的区别，不在于"会不会动"，而在于是否具有以下特征：
\begin{enumerate}
  \item \textbf{有新陈代谢}：能摄取营养、排出废物、进行呼吸。
  \item \textbf{能生长}：从一个小幼苗长成大树，从一个小婴儿长成大人。
  \item \textbf{能繁殖}：制造自己的"副本"。
  \item \textbf{有应激性}：对外界刺激作出反应（含羞草被碰就闭合、向日葵朝向太阳）。
  \item \textbf{能适应和进化}：在漫长的世代中改变自己以适应环境。
\end{enumerate}

有些生物不一定同时表现出所有这些特征——比如骡子能代谢能生长但不能繁殖——但总体上，"活的"和"不活的"之间的界限是清晰的。

\begin{thinkbox}
病毒是活的吗？这是一个经典难题。病毒有遗传物质（DNA或RNA），能繁殖——但它只能在活细胞内繁殖，离开细胞它和一颗灰尘没有区别。大多数生物学家认为病毒"介于生物和非生物之间"。
\end{thinkbox}

\begin{practice}
1. 生物最基本的特征是\_\_\_\_\_。\\
2. （判断）一块岩石也会"长大"（风化后变大），所以它是活的。（\quad）\\
3. 不含羞草被触碰后叶片闭合，体现了生物的\_\_\_\_\_性。\\
4. 病毒为什么不算典型的生物？因为\_\_\_\_\_。\\
5. （判断）植物也能对外界刺激做出反应，如含羞草被碰闭合。（\quad）\\
6. 生物与非生物最根本的区别是生物能进行\_\_\_\_\_。\\
7. 向日葵的花盘随太阳转动，体现了生物的\_\_\_\_\_性。\\
8. 一粒种子种在地里长成一棵大树，体现了生物\_\_\_\_\_的特征。
\end{practice}

\newpage
\section{\textcolor{orange!80!black}{第二课\quad 细胞的分子组成}}

\subsection*{生命的"化学原料"}
组成细胞的主要化学元素（按质量分数）：\textbf{O > C > H > N > Ca > P > S > K > Na > Cl > Mg}。其中O、C、H、N四种元素占细胞干重的96\%以上。

\subsection*{四类生物大分子}

\textbf{糖类（碳水化合物）}：C、H、O。功能：主要能源物质。\\
单糖——葡萄糖C$_6$H$_{12}$O$_6$（最直接的"燃料"）。二糖——蔗糖（甘蔗）、麦芽糖（啤酒）。多糖——淀粉（植物储能）、糖原（动物储能）、纤维素（植物细胞壁）。

\textbf{脂质}：C、H、O（P）。功能：储能、构成膜结构、调节。\\
脂肪（储能）、磷脂（构成细胞膜的双层结构）、固醇（性激素、维生素D的原料）。

\textbf{蛋白质}：C、H、O、N（S）。功能：承担几乎所有生命活动。\\
基本单位——\textbf{氨基酸}（约20种）。多个氨基酸脱水缩合形成肽链——肽键连接。\\
蛋白质的\textbf{结构}决定了其\textbf{功能}：结构不同→形状不同→功能不同（"结构决定功能"是生物学贯穿始终的核心思想）。

\textbf{核酸}：C、H、O、N、P。遗传信息载体。\\
DNA（含有"A、T、G、C"四种碱基）——双螺旋结构，储存在细胞核中。\\
RNA（含有"A、U、G、C"）——参与蛋白质合成的"搬运工"。

\begin{figplaceholder}
此处插入四种生物大分子的结构简图（糖类/脂质/蛋白质/核酸）（见 figure/requirements/bio1-ch1.txt 插图1）
\end{figplaceholder}

\begin{example}
判断：一个氨基酸分子通过肽键连接另一个氨基酸形成了什么呢？

\textbf{解：}两个氨基酸脱水缩合形成\textbf{二肽}。多个氨基酸串联形成\textbf{多肽}，多肽链经过折叠盘绕后才成为有功能的\textbf{蛋白质}。
\end{example}

\begin{practice}
1. 细胞中含量最多的有机物是\_\_\_\_\_；含量最多的无机物是\_\_\_\_\_。\\
2. 蛋白质的基本单位是\_\_\_\_\_；核酸的基本单位是\_\_\_\_\_。\\
3. （判断）葡萄糖的分子式为C$_6$H$_{12}$O$_6$，它属于单糖。（\quad）\\
4. 构成细胞膜的脂质主要是\_\_\_\_\_。\\
5. 糖类的组成元素是\_\_\_\_\_、\_\_\_\_\_、\_\_\_\_\_。\\
6. 人体内最重要的储能物质是\_\_\_\_\_。\\
7. DNA和RNA在碱基上的区别：DNA含\_\_\_\_\_，RNA含\_\_\_\_\_。\\
8. 蛋白质的结构多样性决定了其\_\_\_\_\_多样性。
\end{practice}

\newpage
\section{\textcolor{orange!80!black}{第三课\quad 细胞的结构}}

\subsection*{细胞——生命的基本单位}
\begin{definition}{细胞学说（施莱登和施旺）}
① 一切动植物都由细胞和细胞产物构成；\\
② 细胞是生物体结构和功能的基本单位；\\
③ 新细胞由老细胞分裂产生。
\end{definition}

\subsection*{真核细胞的基本结构}

\textbf{细胞膜}：磷脂双分子层（"海洋"）+ 蛋白质（"冰山"）——流动镶嵌模型。控制物质进出。

\textbf{细胞质与细胞器}：
\begin{itemize}
  \item \textbf{线粒体}：细胞的"发电厂"——有氧呼吸的场所，产生大量ATP。拥有自己的DNA（母系遗传）。
  \item \textbf{叶绿体}（存在植物细胞中）："太阳能电池板"——光合作用的场所。也有自己的DNA。
  \item \textbf{内质网}：蛋白质的加工和运输通道。"粗面"（有核糖体）加工蛋白质；"滑面"合成脂质。
  \item \textbf{高尔基体}：细胞的"快递站"——加工、分拣、运送蛋白质到目的地。
  \item \textbf{核糖体}：蛋白质的"组装车间"——把氨基酸连接成多肽链。
  \item \textbf{溶酶体}：细胞的"垃圾处理厂"——分解废旧细胞器和侵入的细菌。
\end{itemize}

\textbf{细胞核}：细胞的"大脑"——含有遗传物质DNA，指挥细胞的生命活动。核膜（双层膜）、核孔（控制物质进出）、核仁（制造核糖体RNA）。

\subsection*{原核细胞 vs 真核细胞}
\textbf{原核细胞}（细菌、蓝藻）：没有核膜包裹的细胞核——DNA集中在"拟核"区域。没有线粒体、叶绿体等复杂细胞器（只有核糖体）。
\textbf{真核细胞}（动植物、真菌）：有真正的细胞核，有复杂的细胞器。

\begin{practice}
1. 被称为"细胞发电厂"的细胞器是\_\_\_\_\_；被称为"快递站"的是\_\_\_\_\_。\\
2. （判断）原核细胞没有细胞核，所以没有DNA。（\quad）\\
3. 叶绿体和线粒体都拥有自己的\_\_\_\_\_。\\
4. 细胞膜的流动镶嵌模型的主要成分是\_\_\_\_\_和\_\_\_\_\_。\\
5. 原核细胞和真核细胞共有的细胞器是\_\_\_\_\_。\\
6. 植物细胞有而动物细胞没有的细胞器是\_\_\_\_\_和\_\_\_\_\_。\\
7. 细胞核的功能是\_\_\_\_\_。\\
8. 核糖体是合成\_\_\_\_\_的场所。
\end{practice}

\newpage
\section{\textcolor{orange!80!black}{第四课\quad 物质进出细胞}}

\subsection*{扩散与渗透}
\textbf{扩散}：物质从高浓度区域向低浓度区域自发移动——不需要能量。氧气从肺部进入血液、CO$_2$从血液呼出都是扩散。

\textbf{渗透}：水分子通过半透膜从低浓度溶液向高浓度溶液的扩散。植物根毛吸收水分、红细胞在低渗溶液中吸水胀破——都是渗透。

\subsection*{主动运输}
有些物质需要"逆浓度梯度"运输——从低浓度运往高浓度。这需要\textbf{载体蛋白}的帮助和\textbf{能量（ATP）}的消耗。

\textbf{举例：}根细胞从土壤中吸收矿物质离子（土壤中离子浓度远低于细胞内——但根细胞"硬是"把它们泵进来，耗能）。

\subsection*{胞吞与胞吐}
大分子物质进出细胞——通过"囊泡"包裹。白细胞吞噬细菌（胞吞）、神经递质的释放（胞吐）。

\begin{thinkbox}
为什么用0.9\%的生理盐水而不是纯水来输液？因为红细胞在纯水中会吸水胀破——纯水对细胞来说是"极低渗"溶液。0.9\%NaCl溶液的渗透压和血浆一致——细胞既不失水也不吸水。
\end{thinkbox}

\begin{practice}
1. 氧气进入细胞的方式是\_\_\_\_\_，\_\_\_\_\_（需要/不需要）消耗能量。\\
2. 植物根吸收矿物质离子的方式是\_\_\_\_\_运输，因为它\_\_\_\_\_能量。\\
3. （判断）水分子从高浓度溶液向低浓度溶液扩散叫渗透。（\quad）\\
4. 白细胞吞噬细菌的过程叫\_\_\_\_\_。\\
5. 不需要载体蛋白和能量的运输方式是\_\_\_\_\_。\\
6. 葡萄糖进入红细胞的方式是\_\_\_\_\_。\\
7. 胞吞和胞吐依赖的细胞结构是\_\_\_\_\_。\\
8. 成熟植物细胞的原生质层相当于一层\_\_\_\_\_膜。
\end{practice}

\newpage
\section{\textcolor{orange!80!black}{第五课\quad 酶——生命的"催化剂"}}

\subsection*{酶是什么？}
酶是生物体内产生的有催化作用的\textbf{蛋白质}（少数是RNA）。它们能大大降低化学反应所需的活化能——让反应在温和条件下快速进行。

\subsection*{酶的特性}
\textbf{高效性}：酶的催化效率比无机的化学催化剂高$10^7$-$10^{13}$倍。\\
\textbf{专一性}：一种酶只能催化一种（或一类）底物。"一把钥匙开一把锁"。\\
\textbf{敏感性}：酶的活性受\textbf{温度}和\textbf{pH}的影响极大。最适温度：人体酶约37℃，高温使酶变性失活。最适pH：胃蛋白酶约1.5-2，唾液淀粉酶约6.8，通常在中性附近。

\begin{example}
为什么发烧超过40℃人会很危险？除了疾病本身——因为人体内几乎所有酶的最适温度在37℃左右。40℃以上酶开始变性失活——身体的各种代谢反应就会出大问题。
\end{example}

\begin{practice}
1. 绝大多数酶的化学本质是\_\_\_\_\_。\\
2. 酶的特性有\_\_\_\_\_性、\_\_\_\_\_性和\_\_\_\_\_性。\\
3. （判断）高温使酶失活的原因是破坏了酶的空间结构。（\quad）\\
4. 胃蛋白酶的最适pH约为\_\_\_\_\_（填酸性/中性/碱性）。\\
5. 酶催化作用的机理是降低\_\_\_\_\_。\\
6. 唾液淀粉酶的最适pH约为\_\_\_\_\_。\\
7. 过酸或过碱会使酶\_\_\_\_\_。\\
8. 与无机催化剂相比，酶的催化效率更\_\_\_\_\_。
\end{practice}

\newpage
\section{\textcolor{orange!80!black}{第六课\quad 细胞呼吸——"燃烧"食物获取能量}}

\subsection*{有氧呼吸（三步走）}
主要发生在线粒体中。把葡萄糖彻底氧化分解成CO$_2$和H$_2$O，释放大量能量。

\formula{\text{C}_6\text{H}_{12}\text{O}_6 + 6\text{O}_2 \rightarrow 6\text{CO}_2 + 6\text{H}_2\text{O} + \text{能量（ATP）}}

\textbf{三个阶段：}
\begin{enumerate}
  \item \textbf{糖酵解（细胞质基质）}：葡萄糖$\rightarrow$丙酮酸+少量ATP+[H](NADH)。
  \item \textbf{柠檬酸循环（线粒体基质）}：丙酮酸$\rightarrow$CO$_2$+[H](NADH/FADH$_2$)+少量ATP。
  \item \textbf{氧化磷酸化（线粒体内膜）}：[H] + O$_2$ $\rightarrow$ H$_2$O + \textbf{大量ATP}。\\
\end{enumerate}
1分子葡萄糖经有氧呼吸共产生约30-32分子ATP。

\subsection*{无氧呼吸——"急救方案"}
没有氧气时——葡萄糖只能部分分解（糖酵解后不进入线粒体）。\\
\textbf{乳酸发酵}：丙酮酸$\rightarrow$乳酸。人类剧烈运动时肌肉缺氧——产生乳酸（肌肉酸痛的原因）。\\
\textbf{酒精发酵}：丙酮酸$\rightarrow$酒精+CO$_2$。酵母菌在无氧条件下产生酒精（酿酒原理）。

\begin{example}
为什么百米冲刺后腿会又酸又疼？

\textbf{解：}百米冲刺强度极大——氧气供应跟不上肌肉的需求——肌细胞进行\textbf{无氧呼吸}——产生大量\textbf{乳酸}。乳酸在肌肉堆积会刺激神经末梢——产生酸痛感。停下来慢走一会儿、呼吸加快补充氧气——乳酸逐渐被转化为丙酮酸进入有氧循环——酸痛慢慢缓解。
\end{example}

\begin{practice}
1. 有氧呼吸的主要场所是\_\_\_\_\_；无氧呼吸的场所是\_\_\_\_\_。\\
2. （填空）有氧呼吸第三阶段，\_\_\_\_\_和\_\_\_\_\_结合生成H$_2$O，释放大量能量。\\
3. 剧烈运动后肌肉酸痛是因为\_\_\_\_\_在肌肉中积累。\\
4. （判断）有氧呼吸和无氧呼吸第一阶段（糖酵解）是一样的。（\quad）\\
5. 有氧呼吸第二阶段的场所是\_\_\_\_\_。\\
6. 1分子葡萄糖经有氧呼吸可产生约\_\_\_\_\_分子ATP。\\
7. 酵母菌无氧呼吸的产物是\_\_\_\_\_和\_\_\_\_\_。\\
8. 有氧呼吸释放的能量大部分以\_\_\_\_\_形式散失。
\end{practice}

\newpage
\section{\textcolor{orange!80!black}{第七课\quad 光合作用——把光变成"饭"}}

\subsection*{什么是光合作用？}
绿色植物和某些细菌利用光能，把CO$_2$和H$_2$O转化成有机物（主要是葡萄糖），同时放出氧气。

\formula{6\text{CO}_2 + 6\text{H}_2\text{O} \xrightarrow{\text{光}} \text{C}_6\text{H}_{12}\text{O}_6 + 6\text{O}_2}

\subsection*{光反应 vs 暗反应}
\textbf{光反应（类囊体膜上）：}需要光。\\
光$\rightarrow$激发叶绿素$\rightarrow$产生ATP + [H](NADPH) + 释放O$_2$（来自水的光解）。

\textbf{暗反应（叶绿体基质）：}不需要光（但需要光反应提供的ATP和[H]）。\\
CO$_2$ $\rightarrow$ 卡尔文循环$\rightarrow$ 葡萄糖。

\textbf{注意：}暗反应不是晚上才进行——白天也在进行。只要光反应提供了ATP和NADPH，暗反应就能运转。

\subsection*{影响光合作用的因素}
\textbf{光照强度}：在一定范围内，光越强光合越快——但到一定程度后饱和（光饱和点）。\\
\textbf{CO$_2$浓度}：CO$_2$是原料——浓度越高光合越快。\\
\textbf{温度}：影响酶的活性。过高或过低均不利（光合作用涉及一系列酶催化的暗反应）。

\begin{thinkbox}
为什么秋天的树叶会变色？光合作用中的叶绿素在秋季分解后，被"遮盖"的类胡萝卜素（黄/橙）和花青素（红/紫）显露出来。光合作用减弱→树叶"退休"→脱落。我们在自然故事中讲过的——现在是"升级版"的分子解释。
\end{thinkbox}

\begin{practice}
1. 光合作用的场所是\_\_\_\_\_，原料是\_\_\_\_\_和\_\_\_\_\_。\\
2. 光反应的产物有\_\_\_\_\_、\_\_\_\_\_和\_\_\_\_\_。\\
3. （判断）暗反应在夜间进行。（\quad）\\
4. 影响光合作用的四个主要因素：\_\_\_\_\_、\_\_\_\_\_、\_\_\_\_\_、\_\_\_\_\_。\\
5. 光反应发生在叶绿体的\_\_\_\_\_上。\\
6. 暗反应发生在叶绿体的\_\_\_\_\_中。\\
7. 光合作用释放的氧气来自于\_\_\_\_\_。\\
8. 叶绿素主要吸收\_\_\_\_\_光和\_\_\_\_\_光。
\end{practice}

\newpage
\section{\textcolor{orange!80!black}{第八课\quad 细胞增殖——有丝分裂}}

\subsection*{细胞周期}
细胞从一次分裂结束到下一次分裂结束——\textbf{细胞周期}。分为\textbf{间期}和\textbf{分裂期（M期）}。

\textbf{间期}：占细胞周期的大部分时间。DNA复制（每个染色体变成两个姐妹染色单体）、蛋白质合成、能量储备。"准备阶段"。不进行形态分裂。

\textbf{分裂期（有丝分裂）}：核分裂 + 细胞质分裂。分四个阶段（口诀：前期中后末）：

\begin{itemize}
  \item \textbf{前期}：染色质螺旋化成为染色体——每个染色体含两条染色单体。核膜消失。纺锤丝出现。
  \item \textbf{中期}：染色体排列在细胞中央的"赤道板"上。\textbf{染色体形态最清晰、数目最易计数}（观察染色体的最佳时机）。
  \item \textbf{后期}：着丝粒分裂——两条染色单体分离——被纺锤丝拉向两极。染色体数目加倍。
  \item \textbf{末期}：染色体解螺旋成染色质。核膜重新形成。纺锤丝消失。细胞质分裂——一个细胞变成了两个。
\end{itemize}

有丝分裂的意义：产生的两个子细胞\textbf{遗传信息完全相同}（各获得一套与母细胞一样的染色体）。

\begin{figplaceholder}
此处插入有丝分裂四阶段示意图（前期/中/后/末）（见 figure/requirements/bio1-ch1.txt 插图2）
\end{figplaceholder}

\begin{practice}
1. 细胞周期中DNA复制发生在\_\_\_\_\_期。\\
2. 观察染色体形态和数目的最佳时期是\_\_\_\_\_期。\\
3. 有丝分裂后期，着丝粒分裂，染色单体分离成为\_\_\_\_\_。\\
4. （判断）有丝分裂产生的子细胞DNA含量是母细胞的一半。（\quad）\\
5. 细胞周期中分裂间期的主要活动是\_\_\_\_\_。\\
6. 有丝分裂前期染色质螺旋化成为\_\_\_\_\_。\\
7. 有丝分裂末期核膜和\_\_\_\_\_重新形成。\\
8. 有丝分裂的意义是保持子细胞与母细胞\_\_\_\_\_一致。
\end{practice}

\newpage
\section{\textcolor{orange!80!black}{第九课\quad 减数分裂——有性生殖的"特殊分裂"}}

\subsection*{为什么需要减数分裂？}
有性生殖中，精子和卵子各提供一套染色体。如果正常体细胞有$2n$条染色体（二倍体），精子和卵子必须只有$n$条（单倍体）——否则受精后变成$4n$，翻倍了。

\textbf{减数分裂}就是"把染色体数目减半"的特殊分裂。

\subsection*{减数分裂过程}
\textbf{减数第一次分裂}：同源染色体配对（联会）、交叉互换（交换遗传信息）、分离——两个子细胞各得一套染色体（每条含两个染色单体）。

\textbf{减数第二次分裂}：和有丝分裂类似——姐妹染色单体分离——每个子细胞各得一条染色单体。

1个原始生殖细胞 $\rightarrow$ 减数分裂$\rightarrow$ 4个生殖细胞（精子或卵子）。

\subsection*{有丝分裂 vs 减数分裂}
\begin{center}
\begin{tabular}{|c|c|c|}
\hline
& \textbf{有丝分裂} & \textbf{减数分裂} \\
\hline
DNA复制次数 & 1次 & 1次 \\
分裂次数 & 1次 & 2次 \\
子细胞DNA & 与母细胞相同（$2n$） & 减半（$n$） \\
子细胞数目 & 2 & 4 \\
是否发生同源染色体配对 & 否 & 是（前期I联会） \\
\hline
\end{tabular}
\end{center}

\begin{practice}
1. 减数分裂产生的子细胞染色体数目是母细胞的\_\_\_\_\_。\\
2. 同源染色体配对（联会）发生在减数\_\_\_\_\_次分裂的\_\_\_\_\_期。\\
3. （判断）减数分裂过程中DNA复制两次。（\quad）\\
4. 精子和卵子的染色体数目是体细胞的\_\_\_\_\_。\\
5. 同源染色体联会发生在减数\_\_\_\_\_次分裂的\_\_\_\_\_期。\\
6. 减数第二次分裂姐妹染色单体分离发生在\_\_\_\_\_期。\\
7. 交叉互换发生在减数第一次分裂的\_\_\_\_\_期。\\
8. 一个精原细胞经减数分裂能产生\_\_\_\_\_个精子。
\end{practice}

\newpage
\section{\textcolor{orange!80!black}{第十课\quad 细胞分化、衰老与凋亡}}

\subsection*{分化——同源细胞走向不同命运}
一个受精卵$\rightarrow$神经细胞、肌肉细胞、血细胞……这些细胞都有完全相同的DNA——但为什么形态和功能差异巨大？因为\textbf{基因的选择性表达}——每个细胞只"打开"它需要的基因。

\textbf{干细胞}：尚未分化、能分化成多种细胞类型的细胞。胚胎干细胞（全能性最强）和成体干细胞（如骨髓造血干细胞——能分化成各种血细胞）。干细胞研究是再生医学的核心。

\subsection*{细胞的衰老}
细胞不会永远分裂下去——正常细胞分裂次数是有限的（Hayflick极限，约50-60次）。衰老的细胞代谢减慢、酶活性降低、线粒体功能下降。

\subsection*{细胞的凋亡——"程序性死亡"}
不是意外死亡——是基因控制的、「主动」结束生命的过程。蝌蚪变青蛙时尾巴消失（尾部的细胞凋亡了）、人类胚胎发育时手指之间的蹼消失了（细胞凋亡"挖出"了手指缝）。

\textbf{区别：}坏死是"外伤致死"（引起炎症）；凋亡是"按计划赴死"（不引起炎症）。

\begin{thinkbox}
癌细胞的特点是"永生"——不受控制地无限分裂。正常细胞分裂到一定次数就"决定"了（衰老/凋亡）。癌细胞把控制分裂的机制"搞坏了"——不停地分裂、抢夺营养、侵犯周围组织。这就是癌症最可怕的地方——它不是"外界入侵"，而是自身的细胞"叛变了"。
\end{thinkbox}

\begin{practice}
1. 不同组织细胞形态和功能不同，是因为\_\_\_\_\_表达不同。\\
2. （判断）干细胞具有全能性，可以分化成多种细胞。（\quad）\\
3. 蝌蚪尾巴消失的原因是细胞\_\_\_\_\_。\\
4. 细胞凋亡和坏死的区别是\_\_\_\_\_。\\
5. 受精卵能分化成各种细胞，说明它具有\_\_\_\_\_。\\
6. 正常细胞分裂次数有限，这被称为\_\_\_\_\_极限。\\
7. 癌细胞的特点是\_\_\_\_\_。\\
8. 细胞坏死会引起\_\_\_\_\_（填"炎症"或"不炎症"）。
\end{practice}

\vspace{0.5cm}
\begin{center}\rule{0.7\textwidth}{1pt}\vspace{0.4cm}
{\large\textcolor{blue!70!black}{\textbf{—— 生物1·分子与细胞 完 ——}}}
\end{center}

\newpage
\section*{\textcolor{defblue}{参考答案}}

\subsection*{第一课}
1. 新陈代谢\\
2. （×）\\
3. 应激\\
4. 病毒不能独立进行新陈代谢，必须寄生在活细胞内才能繁殖\\
5. （√）\\
6. 新陈代谢\\
7. 应激\\
8. 生长

\subsection*{第二课}
1. 蛋白质、水\\
2. 氨基酸、核苷酸\\
3. （√）\\
4. 磷脂\\
5. C、H、O\\
6. 脂肪\\
7. T（胸腺嘧啶）、U（尿嘧啶）\\
8. 功能

\subsection*{第三课}
1. 线粒体、高尔基体\\
2. （×）原核细胞有DNA（拟核中）\\
3. DNA\\
4. 磷脂（磷脂双分子层）、蛋白质\\
5. 核糖体\\
6. 叶绿体、液泡\\
7. 含有遗传物质DNA，指挥细胞生命活动\\
8. 蛋白质

\subsection*{第四课}
1. 扩散（自由扩散）、不需要\\
2. 主动、消耗\\
3. （×）水从低浓度溶液向高浓度溶液扩散叫渗透\\
4. 胞吞\\
5. 自由扩散\\
6. 协助扩散（易化扩散）\\
7. 囊泡（细胞膜）\\
8. 半透

\subsection*{第五课}
1. 蛋白质（少数是RNA）\\
2. 高效、专一、敏感\\
3. （√）\\
4. 酸性\\
5. 活化能\\
6. 6.8（接近中性）\\
7. 变性失活\\
8. 高

\subsection*{第六课}
1. 线粒体、细胞质基质\\
2. [H]（NADH）、O$_2$\\
3. 乳酸\\
4. （√）\\
5. 线粒体基质\\
6. 30--32\\
7. 酒精（乙醇）、CO$_2$\\
8. 热能

\subsection*{第七课}
1. 叶绿体、CO$_2$、H$_2$O\\
2. ATP、[H]（NADPH）、O$_2$\\
3. （×）暗反应不需要光，但白天也能进行\\
4. 光照强度、CO$_2$浓度、温度、水分\\
5. 类囊体（膜）\\
6. 基质\\
7. 水（H$_2$O的光解）\\
8. 红、蓝紫

\subsection*{第八课}
1. 间\\
2. 中\\
3. 染色体\\
4. （×）DNA含量与母细胞相同\\
5. DNA复制和蛋白质合成\\
6. 染色体\\
7. 核仁（或纺锤丝）\\
8. 遗传信息

\subsection*{第九课}
1. 一半（$n$）\\
2. 第一次分裂的前（I前期）\\
3. （×）DNA复制一次\\
4. 一半（$n$）\\
5. 第一次分裂的前（I前期）\\
6. 后（II后期）\\
7. 前（I前期）\\
8. 4

\subsection*{第十课}
1. 基因的选择性\\
2. （√）\\
3. 凋亡（程序性死亡）\\
4. 凋亡是基因控制的程序性死亡（不引起炎症）；坏死是意外损伤（引起炎症）\\
5. 全能性\\
6. Hayflick\\
7. 不受控制地无限分裂（永生）\\
8. 炎症

\end{document}
